\section{A Four-Axis Taxonomy}
\label{sec:taxonomy}

We organize the landscape along four axes (\Cref{fig:taxonomy}).

\begin{figure}[t]
\centering
\begin{tikzpicture}[
    axis/.style={draw, rounded corners, minimum width=2.8cm, minimum height=0.7cm,
                 font=\small\bfseries, align=center, fill=blue!8},
    sub/.style={font=\scriptsize, align=left, text width=2.6cm},
    >=Stealth
]
% Four axes arranged in a 2x2 grid
\node[axis] (s) at (0, 0) {Structure};
\node[axis] (r) at (4, 0) {Retrieval};
\node[axis] (l) at (0, -2.2) {Lifecycle};
\node[axis] (p) at (4, -2.2) {Preference};

% Sub-labels
\node[sub, below=2pt of s.south west, anchor=north west]
    {Working, Episodic,\\Semantic, Procedural,\\Implicit};
\node[sub, below=2pt of r.south west, anchor=north west]
    {Na\"ive RAG,\\Advanced RAG,\\Modular RAG};
\node[sub, below=2pt of l.south west, anchor=north west]
    {Formation,\\Consolidation,\\Forgetting, Transform.};
\node[sub, below=2pt of p.south west, anchor=north west]
    {None, Extracted,\\Opinion-tracked,\\Emergent};

% Orthogonality indicators
\draw[<->, gray, dashed] (s.east) -- (r.west);
\draw[<->, gray, dashed] (s.south) -- (l.north);
\draw[<->, gray, dashed] (r.south) -- (p.north);
\draw[<->, gray, dashed] (l.east) -- (p.west);
\end{tikzpicture}
\caption{Four-axis taxonomy for AI agent memory systems. Dashed arrows
indicate orthogonality: a system's position on one axis places no
constraint on the others.}
\label{fig:taxonomy}
\end{figure}

\subsection{Axis 1: Memory Structure}

Following CoALA~\cite{sumers2024coala}, we classify memory modules as:

\begin{itemize}[leftmargin=*, nosep]
    \item \textbf{Working memory:} In-context information the LLM
    processes directly. Managed by the agent, not the memory system.
    \item \textbf{Episodic memory:} Raw experiences preserved with
    full context (who, when, where, about what).
    \item \textbf{Semantic memory:} Distilled, decontextualized
    knowledge extracted from episodes.
    \item \textbf{Procedural memory:} Executable skills and
    behavioral patterns.
    \item \textbf{Implicit memory:} Preferences and habits that
    emerge from accumulated interaction without explicit extraction.
    Not present in CoALA; we propose it as an extension.
\end{itemize}

\Cref{tab:memory-type-dist} summarizes how these types are distributed
across the 112 systems in our survey. Note that systems may implement
more than one type, so percentages sum to greater than 100\%.

\begin{table}[t]
    \centering
    \caption{Memory type distribution across 112 surveyed systems.}
    \label{tab:memory-type-dist}
    \begin{tabular}{lcc}
        \toprule
        \textbf{Memory Type} & \textbf{Systems} & \textbf{\%} \\
        \midrule
        Working      & 40 & 36 \\
        Episodic     & 64 & 57 \\
        Semantic     & 58 & 52 \\
        Procedural   &  8 &  7 \\
        Implicit     & 35 & 31 \\
        \bottomrule
    \end{tabular}
\end{table}

Most systems implement only one or two memory types.
Episodic and semantic memory dominate, appearing in a majority of
systems, while procedural and implicit memory remain rare---present in
only 7\% and 31\% of systems, respectively.
This asymmetry reflects the field's current bias toward
\emph{storage} (retaining raw experiences and extracted facts) over
\emph{learning} (acquiring reusable skills or evolving preferences from
accumulated interaction).

\subsection{Axis 2: Retrieval Architecture}

Following Gao~\etal~\cite{gao2023rag}, we classify retrieval as:

\begin{itemize}[leftmargin=*, nosep]
    \item \textbf{Naive RAG:} Single retrieval signal (vector or
    keyword), direct injection. Examples: LangChain, Engram.
    \item \textbf{Advanced RAG:} Multiple signals or reranking.
    Examples: Generative Agents (recency $\times$ importance $\times$
    relevance), Mem0 (vector + graph + priority).
    \item \textbf{Modular RAG:} Composable pipeline with independent
    retrieval, fusion, and reranking stages. Examples: Zep/Graphiti
    (cosine + BM25 + graph + RRF + MMR), SYNAPSE (multi-signal with
    activation-based scoring).
\end{itemize}

\subsection{Axis 3: Memory Lifecycle}

Following Zhang~\etal~\cite{zhang2024survey} and
Hu~\etal~\cite{hu2025memory}:

\begin{itemize}[leftmargin=*, nosep]
    \item \textbf{Formation:} How memories enter the system (direct
    storage, LLM extraction, agent-directed).
    \item \textbf{Consolidation:} How episodic memories transform into
    semantic knowledge (CLS replay, LLM reflection, sleep-time
    processing).
    \item \textbf{Forgetting:} How memories weaken or disappear
    (no forgetting, single-curve decay, dual-strength decay, RL-learned).
    \item \textbf{Transformation:} How the memory store itself
    restructures (deduplication, contradiction resolution, pruning).
\end{itemize}

\subsection{Axis 4: Preference Emergence}

We propose this as a new axis not covered by existing surveys:

\begin{itemize}[leftmargin=*, nosep]
    \item \textbf{No preferences:} Most systems.
    \item \textbf{LLM-extracted profiles:} Mem0, Supermemory, Memobase.
    One-shot extraction, brittle.
    \item \textbf{Explicit opinion tracking:} Hindsight. Confidence-scored
    beliefs that evolve.
    \item \textbf{Emergent crystallization:} Preferences crystallize
    from accumulated impressions without per-interaction LLM involvement.
    Rare; only 4.5\% of surveyed systems.
\end{itemize}

\medskip

These four axes are orthogonal: a system's position on one axis places
no constraint on the others. For instance, HippoRAG employs modular RAG
retrieval but implements no lifecycle operations, while MemoryBank uses
naive single-signal retrieval yet provides sophisticated lifecycle
management through Ebbinghaus forgetting curves. The axes thus capture
genuinely independent dimensions of the design space, and a complete
characterization of any memory system requires specifying its position
on all four.
